\section{System Design}
\label{sec:design}

\subsection{Architecture}

ALS is designed as a \emph{centralized audit collector} following the same
supporting-service pattern as the PRE Server: an independent Rust binary with
a single well-defined responsibility, reusing DecMed's existing IOTA and IPFS
infrastructure. Four event sources feed into ALS: Patient Client, Hospital
Client, Ministry Client, and PRE Server.

ALS comprises four internal components:

\begin{itemize}
    \item \textbf{Audit Receiver}: an Axum HTTP endpoint
    (\texttt{POST /api/events}) that accepts encrypted-and-signed event
    payloads from event sources.

    \item \textbf{Internal Queue}: a bounded async channel (capacity 10{,}000)
    that decouples event reception from processing, preventing back-pressure
    on event sources.

    \item \textbf{Audit Logger}: an async worker that dequeues events,
    verifies signatures, forms structured \texttt{AuditRecord}s, applies
    hash-chaining, and appends records to the active log file.

    \item \textbf{Log Rotation \& Archival Worker}: a timer-driven worker that
    seals the active log file, uploads it to IPFS Cluster with pinning, and
    commits rotation metadata to IOTA.
\end{itemize}

\subsection{Event Encapsulation and Non-Repudiation}

Non-repudiation requires that the origin of each audit event be
cryptographically bound to the sending component. Each event source generates
an ephemeral Ed25519 key pair at initialization. When an auditable operation
completes, the source constructs an \texttt{AuditEvent}, serializes it, and
signs it with its Ed25519 private key. The signed payload is then encrypted
under ALS's ECDH P-256 public key using ECIES (AES-256-GCM for symmetric
encryption, SHA-256 as the KDF), yielding an \texttt{EncryptedSignedEvent}.

ALS decrypts the payload, verifies the Ed25519 signature against the
registered public key of the source component, and only then enqueues the
event. Events with invalid or absent signatures are rejected and themselves
logged as EV12 (policy violation) entries.

The asynchronous delivery model means that ALS events do not add latency
to the critical path of clinical operations. Each event source integrates an
\texttt{ALSClient} backed by a local \texttt{AlsQueue}; a retry worker
re-attempts delivery with exponential backoff if ALS is temporarily
unavailable, so events survive transient ALS downtime without being dropped.

\subsection{Audit Record Schema and Hash-Chaining}

Each verified event is materialized into an \texttt{AuditRecord}:

\begin{itemize}
    \item \texttt{record\_id}: UUID v7 (monotone, time-ordered).
    \item \texttt{timestamp}: RFC~3339 UTC, set at materialization time by ALS.
    \item \texttt{prev\_record\_hash}: SHA-256 hex digest of the preceding
    record's serialization, or \texttt{null} for the first record in a log
    file.
    \item Common fields: \texttt{source\_component}, \texttt{actor\_id},
    \texttt{actor\_type}, \texttt{action} (EV1--EV13), \texttt{target\_object\_id},
    \texttt{target\_type}, \texttt{outcome}, \texttt{event\_signature}.
    \item \texttt{details}: a tagged enum (\texttt{AuditEventDetails})
    carrying event-specific evidence attributes.
\end{itemize}

Records are written in JSON-Lines format to an append-only local log file.
The \texttt{prev\_record\_hash} field of record $n$ contains
$\text{SHA-256}(\text{serialize}(r_{n-1}))$. Any modification, insertion, or
deletion of a record breaks the chain at that position and is detected
during pre-rotation verification before the file is sealed and uploaded.

\subsection{Log Rotation and Long-Term Retention}

IPFS's content-addressed storage prohibits in-place append after a file is
published; each upload produces a new CID. ALS therefore buffers records
locally and applies a time-based rotation strategy:

\begin{enumerate}
    \item At the configured rotation interval, the Audit Logger seals the
    active log file (verifying the hash chain) and opens a new one.

    \item The sealed file is uploaded to IPFS Cluster. The cluster applies
    \emph{pinning} to mark the file as permanent, preventing garbage
    collection.

    \item ALS submits a Move transaction via the existing Gas Station Server
    to call \texttt{als::audit\_log::commit\_rotation}. The resulting on-chain
    object stores the CID, file hash (SHA-256), record count, and time period
    as a \emph{frozen} (immutable) object.
\end{enumerate}

This dual-layer strategy eliminates both failure modes: IPFS pinning prevents
GC removal of the log archive; the on-chain frozen object prevents IOTA node
pruning from erasing the metadata anchor. The IOTA frozen object cannot be
modified by any actor, including the ALS operator, after the transaction is
finalized.